\section{Design Considerations}
\label{chap:design_decisions}

\textbullet\ \textbf{ I$^2$C Master-Only Mode:}

One of the initial design decisions was to focus the development of the 
\\STM32F405's  I$^2$C interface exclusively on master mode. This is because, in the 
educational context in which it will be used, there is no need to use a microcontroller 
as an  I$^2$C slave. Furthermore, the LIS3DH sensor does not have a master mode, so 
this decision defines the operational context in which the STM32F405 communicates 
with the LIS3DH accelerometer. This decision simplifies the scope of the emulation 
while maintaining full functional reliability for the intended use case.

\textbullet\ \textbf{Timing Register Abstraction:}

QEMU implements the  I$^2$C protocol through functions that abstract the different 
events of the  I$^2$C protocol (Code~\ref{code:qemu_i2c_fun}), without requiring cycle-level clock generation. This 
abstraction significantly reduces implementation complexity and preserves functional correctness for firmware testing. 
To emulate the STM32F405 I$^2$C protocol, this means that it is not necessary 
to implement the timing logic controlled by the I2C\_CCR and I2C\_TRISE registers.

\lstinputlisting[
  language=C,
  caption={\texttt{qemu i2c functions}},
  label={code:qemu_i2c_fun}
]{Resources/Code/i2c_functions.c}


\textbullet\ \textbf{Netduino Plus 2 as a Development Board for the STM32F405 Microcontroller:}

Unlike other emulators, QEMU fully emulates a hardware platform; it doesn't 
offer the option to emulate only MCU models independently. This limitation 
created the need to find or implement a development board model in QEMU that 
used the STM32F405 microcontroller. Fortunately, QEMU has a pre-implemented 
minimalist model of the Netduino Plus 2 electronic development board, based on 
the STM32F405RG microcontroller, which has saved users from having to develop 
a new board to implement the microcontroller.

\textbullet\ \textbf{DRDY-Only Interrupt Implementation:}

As explained during the analysis of LIS3DH in section~\ref{chap:analysis_lis3dh}, this accelerometer 
has several interrupts. Of all the interrupts it offers, only the data-ready (DRDY) 
interrupt logic could be implemented. While the goal was to implement the logic 
for all interrupts, the development began with the DRDY interrupt logic due to its 
simplicity and because it represents the most common interrupt mode in embedded 
applications. Unfortunately, due to time constraints, the logic for the remaining 
interrupts could not be implemented. These omissions do not affect the basic 
functionality required for typical STM32 firmware development scenarios, where DRDY 
polling or simple data ready interrupts are sufficient.

\textbullet\ \textbf{QOM as interface:}

Conceived as a temporary solution until its full implementation in MICROLAB, 
the LIS3DH accelerometer model has been designed to allow runtime access to 
acceleration values for each axis. This has been achieved by defining QOM properties 
in the model. As explained during the QEMU analysis, the QOM API allows the 
creation of dynamic properties that can be modified at runtime. Thanks to these 
properties, the developer can inject known acceleration values using commands, such 
as those described in the Figure~\ref{fig:qom_exampe}, to test the accelerometer's functionality with known 
values.

\begin{figure}[H]
    \centering
    \begin{verbatim}
        (qemu) qom-set /machine/soc/i2c[0]/i2c1/child[0] accel-x 500 
        (qemu) qom-set /machine/soc/i2c[0]/i2c1/child[0] accel-y -300 
        (qemu) qom-set /machine/soc/i2c[0]/i2c1/child[0] accel-z 980
        
        (qemu) qom-get /machine/soc/i2c[0]/i2c1/child[0] accel-x 
        500
        (qemu) qom-get /machine/soc/i2c[0]/i2c1/child[0] accel-y 
        -300
        (qemu) qom-get /machine/soc/i2c[0]/i2c1/child[0] accel-z
        980  
    \end{verbatim}
    \caption{access to qom properties}
    \label{fig:qom_exampe}
\end{figure}

